% chapters/chapter05.tex
\chapter{Che lavoro fa Mark?}
«Bene, allora lo inoltriamo.»\\
«Sì... grazie.» La voce di Caterina \`e un sospiro. Con un dito Laura chiude la procedura e con una mano il computer portatile. «Questa \`e fatta, per\`o non \`e sufficiente. Bisognerà sbloccare la situazione, altrimenti vedrai che te lo negheranno ancora.»\\
«Grazie Laura, vediamo cosa succede allora.»\\
«Sì, vediamo. Adesso per\`o ceniamo.» Si gira verso la sorella che  tra le dita dei piedi ha infilato  i cappucci dei pennarelli. Trattiene un moto di impazienza, li sfila  e  tappa i colori:  «Vale vai a lavarti le mani... e i piedi!»\\
«Deve venire anche Caterina per\`o!»\\
«Sì, vengo anche  io. E Cerbero? Non si lava le zampe?» Laura sorride all'amica, Caterina \`e pi\`u rilassata e si vede. A volte Caterina si agita per un nonnulla e si tranquillizza per ancora meno.\\
Laura, ora sola in cucina, aggiunge piatto, posate e bicchiere per l'amica. Guarda il computer sul tavolino, la aspetta un dopocena in una rete di connessioni impalpabili che Gibson avrebbe definito cyberspazio.

\par\medskip
Rocky \`e seduto ai piedi di Cate. Gli occhi sono attenti, lui ha gi\`a dimenticato ogni boccone trafugato e briciola caduta. Conta solo il presente.  Segue ogni movimento della cena che, suo malgrado, \`e ormai giunta al termine.
«Bene», dice Laura mentre si alza per sparecchiare, ma non fa in tempo a voltarsi che la sorella  scappa sul tappeto per finire i suoi disegni. Un  «Torna subito qui» va a vuoto ed \`e Caterina che l'aiuta con i piatti.\\
«Dì quello che hai in mente, sono tre minuti che asciughi lo stesso bicchiere.»\\
Lo ripone nella credenza e si lascia andare: «Lo sapevo che il lavoro di Mark prima o poi mi avrebbe ostacolata.»\\
«Cosa ti preoccupa di pi\`u, il corso o Mark?»\\
«Non lo so, per\`o non sono contenta. Sarebbe un problema anche spiegare che non ho avuto il visto...»\\
«Ma nel tuo fascicolo alla Pet-Micro-Robots risulta il tuo passato di attivista?»\\
«In parte... certo non gli ho raccontato dell'arresto.»\\
Laura spalanca gli occhi e per un attimo lascia scorrere l'acqua senza fermarla: «Cosa? Ti hanno arrestata?»\\
«Sì, ci hanno portato alla centrale della Polizia Locale, ma non siamo stati denunciati, solo... come si dice, ci hanno preso le generalità.»\\
«In che occasione?» Laura deve approfondire perch\'e forse la spiegazione \`e vicina.\\
«Cio\`e cosa stavamo facendo? Niente, una manifestazione contro i nuovi impianti di micro reattori in aree T1.»\\
«Tre anni fa? La protesta contro la SAGA-M?»\\
Caterina si stringe nell'imbarazzo: «All'epoca il nome non mi diceva nulla... ma sì, in effetti erano della SAGA-M.» Appoggia il piatto sulla tavola. «Quindi? \`E per quello?» le mani le tremano.
Laura chiude l'acqua del lavello e con le mani ancora bagnate stringe le sue: «Aspetta, stai tranquilla...»\\
«Sarebbe il colmo, non posso andare al mio corso perch\'e non vado a genio alla compagnia del mio ragazzo.»\\
«Cate...»\\
«...scusami, non ce l'ho con te. Adesso mi calmo.»\\
«Cate, non \`e ancora tutto perduto.»\\
«Ma Laura: potrei perdere il mio lavoro. Tu sai che avventura abbiamo vissuto per averlo\footnote{Gli eventi cui Caterina fa riferimento sono narrati in
\emph{CNOT -- Fuga dal computer quantistico}, attualmente disponibile
in versione digitale su GitHub: \url{https://github.com/francescosisini/Cnot-Franchise}.}.»\\
Uno sguardo, un sorriso e si abbracciano: «Appunto, ce la siamo cavata, no?»\\
Dal tappeto Valentina si gira per guardarle: «Dovete compiere un'altra missione? Posso venire anche io?»\\
Valentina ha le palpebre pesanti. «Vai a lavare i denti che \`e quasi ora...»\\
«Non voglio andare a letto» protesta. «Dai, dimmi cosa dovete fare.»\\
«Te lo dico domani, adesso corri in bagno.»\\
Caterina si toglie il grembiule e sistema le maniche della camicetta. Espira forte, come ha imparato dal maestro di yoga: «Sì, vado anche io. Fra poco torna a casa anche Alice.»\\
Laura le accarezza i capelli:  «Un episodio di tre anni fa che risulta ancora visibile \`e... non dico strano ma, particolare.»\\
«Quindi non \`e quello?»\\
«Potrebbe, ma non dovrebbe. La metto a letto poi cerco di capire, ok? Tu vai a casa da Alice, se mi serve qualcosa ti chiamo.»\\

«Ma Laura, non volevo disturbarti ancora.»\\
«Non preoccuparti, \`e solo allenamento per il prossimo esame.»\\


\par\medskip
Casa di Laura \`e un po' in periferia, una zona scura, con poche luci artificiali. Caterina  cerca l'Orsa Maggiore. Valentina l'aveva rappresentata nel cielo del suo disegno e le aveva  raccontato una leggenda a riguardo. I satelliti artificiali per\`o percorrendo le loro orbite alterano le relazioni tra le stelle fisse; ogni epoca ha le sue difficoltà. Finalmente riconosce la Polare e poi il Carro, ma un satellite ne attraversa lentamente la coda. Ecco cosa le aveva detto: «\`E un segno di mutamento,   annuncia una cosa grossa!». Un brivido o una brezza leggera le sfiora le spalle: «Fa fresco» le sfugge.\\
«Mi sa che hai dimenticato la giacca in casa. Te la prendo» risponde prontamente Laura.

\par\medskip
Caterina rimane sola. Poco lontano vede un cane camminare da solo, forse \`e randagio, o forse si \`e perso. Vuole dirlo a Laura ma un movimento rapido in cielo la distrae; alza lo sguardo di nuovo verso le stelle: \`e solo un drone come tanti.\\
«Eccolo, Cate.»\\
«Grazie.»\\
«Allora, mi aggiorni?»\\
«Certo.»

\par\medskip
Caterina si dirige verso casa, da Alice, mentre Laura chiude la vetrata e blocca la porta di ingresso.

\par\medskip
«Andiamo a letto?»\\
La voce di Valentina la raggiunge. La sente forte, come un rimprovero.

\par\medskip
I futon sono gi\`a stesi per terra. Valentina ha in mano un tankōbon di \textit{Buonanotte, Punpun} che apparteneva alla loro mamma.  Laura ha fermato alcune pagine con le graffette e lo legge ogni sera insieme a lei prima di dormire.\\
«Hai proprio sonno!»\\
«No, \`e che voglio leggere Punpun. Dai, vieni.»\\
«Ma... ok, per\`o solo un capitolo, perch\'e poi devo fare una cosa per Cate.»\\
«Cosa devi fare? La vostra avventura?»\\
«Sì, proprio quello. Domani ti racconto tutto.»


\par\medskip
Leggono cinque pagine, poi Vale sbadiglia:
«Hai sonno? Finiamo il capitolo?»\\
«\`E lo stesso, adesso dormo. Tu cosa fai?»
«Devo fare quella cosa per Caterina te l'ho detto.»\\
«Posso venire con te?»\\
«Ma non hai detto che hai sonno?»\\
«Sì, ma voglio vedere che...» Non finisce la frase. Laura \`e abituata a vederla crollare di sera.\\
Lascia la porta della loro stanza socchiusa. Prende la caraffa di t\`e freddo dal frigo e se ne versa  un bicchiere. Il computer \`e gi\`a sul tavolo, Punpun \`e tornato nella libreria.
%Un DPL: dispositivo di programmazione con kernel Linux.
«Vediamo la situazione»:
\begin{terminalbox}
kali@dpl:~$ dig +short polizialocale.ith56.it
203.0.113.42
203.0.113.57
203.0.113.81
203.0.113.96

kali@dpl:~$
\end{terminalbox}

\par\medskip
%I dati di Cate devono svanire dal server giusto un attimo prima che il DHS, via TECS, chieda l’accesso. Poi torneranno al loro posto  come se nulla fosse.

\par\medskip
\begin{center}
  \includegraphics[width=.9\linewidth]{rete senza sfondo.png}
\end{center}

\par\medskip
«128 minuti per mappare l'infrastruttura e segnare i punti deboli. Prima una perlustrazione, vediamo le porte vintage. Se va bene trovo anche una SQL injection... Ce la possiamo fare, Rocky.»

\par\medskip
Allunga la mano, d’istinto, verso la cuccia.

\par\medskip
«Rocky? Dove sei, Rocky?»
